\chapter{Mémoire torsionnelle et conditions de rebond dans Einstein--Cartan}
\label{ch:rebote-ec}

\section{Transport de la loi de Perron vers la densité spectrale de spin}

La réalisation cosmologique de Perron construit deux séquences sur un
même chemin fermé :

\begin{equation}
 \frac{a_n}{a_0}=\varphi^{n/3},
 \qquad
 M_n=\varphi^{-2n}.
 \label{eq:rebote-perron-pair}
\end{equation}

L'élimination du niveau \(n\) produit l'identité exacte

\begin{equation}
 \boxed{M_n=\left(\frac{a_n}{a_0}\right)^{-6}.}
 \label{eq:rebote-memory-six}
\end{equation}

L'exposant possède une généalogie déterminée : la dimension spatiale
apporte le facteur \(3\) et le caractère quadratique de la mémoire stable
apporte le facteur \(2\). La puissance \(6=2\cdot3\) exprime ainsi la dilution
d'une densité quadratique associée à une charge extensive de spin.

\begin{definition}[Transducteur torsionnel]
Une amplitude positive \(s_*^2\) étant fixée dans une section initiale, le transducteur
torsionnel est
\begin{equation}
 \mathcal T_{\rm EC}:M_n\longmapsto
 s_n^2=s_*^2M_n
 =s_*^2\left(\frac{a_n}{a_0}\right)^{-6}.
 \label{eq:rebote-transductor}
\end{equation}
Son domaine est le mode stable de Perron et son codomaine est le secteur
quadratique de spin d'une réalisation d'Einstein--Cartan.
\end{definition}

L'équation \cref{eq:rebote-transductor} conserve l'exposant et la
positivité. L'amplitude \(s_*^2\), le signe avec lequel la contorsion intervient dans
l'équation gravitationnelle et la normalisation dimensionnelle appartiennent au
réalisateur physique. Cette séparation situe exactement ce qu'apporte le
tronc HMT et ce que doit apporter la dynamique de Cartan.

\section{Équations homogènes et convention de densité}

Pour exposer le mécanisme structurel, nous adoptons temporairement des unités
\(c=1\) et une géométrie FRW spatialement plate. Soit

\begin{equation}
 \rho_m(a)=\rho_0a^{-3(1+w)},
 \qquad
 \rho_s(a)=s_0a^{-6},
 \label{eq:rebote-densities}
\end{equation}

avec \(\rho_0>0\), \(s_0>0\) et \(w\) constant. La contribution de
contorsion intervient avec un signe répulsif dans l'équation effective :

\begin{equation}
 H^2=\frac{8\pi G}{3}\bigl(\rho_m-\rho_s\bigr).
 \label{eq:rebote-friedmann}
\end{equation}

Si l'on emploie des densités d'énergie, le membre droit est multiplié par
\(c^{-2}\) ; si l'on emploie des densités de masse, il conserve la forme
\cref{eq:rebote-friedmann}. Le document maintient une seule convention dans
chaque calcul et restaure la carte dimensionnelle à la fin.

L'équation de Raychaudhuri correspondante est

\begin{equation}
 \dot H=-4\pi G\bigl[(1+w)\rho_m-2\rho_s\bigr].
 \label{eq:rebote-raychaudhuri}
\end{equation}

La deuxième contribution se comporte comme un fluide rigide d'équation
d'état \(w_s=1\) et de densité effective négative. Cette forme est la publication
homogène de la répulsion torsionnelle à grande densité.

\section{Existence et unicité du rayon de rebond}

\begin{theorem}[Rebond régulier pour \(w<1\)]
Supposons \(\rho_0>0\), \(s_0>0\) et \(w<1\). Il existe alors un unique
rayon positif \(a_b\) tel que \(H(a_b)=0\), donné par
\begin{equation}
 \boxed{
 a_b^{3(1-w)}=\frac{s_0}{\rho_0}.}
 \label{eq:rebote-radius}
\end{equation}
En ce point
\begin{equation}
 \boxed{
 \dot H_b=4\pi G(1-w)\rho_b>0,}
 \qquad
 \rho_b:=\rho_m(a_b)=\rho_s(a_b),
 \label{eq:rebote-positive-derivative}
\end{equation}
de sorte que la branche contractante se prolonge en une branche expansive régulière.
\end{theorem}

\begin{proof}
La condition \(H=0\) équivaut à
\(\rho_0a^{-3(1+w)}=s_0a^{-6}\). Multiplier par \(a^6\) produit
\cref{eq:rebote-radius}. L'exposant \(3(1-w)\) est positif ; la
puissance définit donc une bijection de \(\RR_{>0}\) sur lui-même et le rayon est
unique. En substituant \(\rho_m=\rho_s=\rho_b\) dans
\cref{eq:rebote-raychaudhuri},
\[
 \dot H_b=-4\pi G[(1+w)-2]\rho_b
 =4\pi G(1-w)\rho_b>0.
\]
La positivité transforme le zéro simple de \(H\) en un minimum de \(a(t)\).
\end{proof}

Le développement local autour de \(t_b\) rend la régularité visible :

\begin{equation}
 a(t)=a_b\left[
 1+2\pi G(1-w)\rho_b(t-t_b)^2
 +O\!\left((t-t_b)^3\right)
 \right].
 \label{eq:rebote-local-expansion}
\end{equation}

La loi \(a^{-6}\) apporte la puissance qui domine le secteur matériel lorsque
\(w<1\) et que \(a\) diminue. Le signe torsionnel et l'amplitude positive font
que cette dominance clôt la contraction avant d'atteindre \(a=0\).

\section{Cas limites de la dynamique torsionnelle et obstructions à la régularité du rebond}

La valeur \(w=1\) produit le même exposant \(a^{-6}\) dans les deux termes.
La différence conserve alors un signe fixe déterminé par
\(\rho_0-s_0\) ; l'égalité ponctuelle cesse de sélectionner un rayon isolé.
Pour \(w>1\), la densité matérielle croît plus vite que le terme
torsionnel pendant la contraction et le mécanisme précédent ne garantit pas de
rebond. Ces cas fournissent des contrôles négatifs internes : la clôture
dépend de l'inégalité des exposants et non d'un choix rétrospectif
du seuil.

La courbure spatiale \(k\), une constante cosmologique et les anisotropies
de Bianchi ajoutent des termes à \cref{eq:rebote-friedmann}. Leur effet local
se contrôle sans présupposer la position du nouveau zéro. Définissons
\[
 F_0(a):=\frac{8\pi G}{3}\bigl(\rho_m(a)-\rho_s(a)\bigr).
\]
L'équation perturbée prend la forme
\(H^2=F_R(a):=F_0(a)+R(a)\).

\begin{theorem}[Persistance \(C^1\) du rebond transversal]
\label{thm:rebote-c1-persistence}
Soit \(a_b\) le rayon de \cref{eq:rebote-radius}. Il existe
\(0<a_-<a_b<a_+\) tels que, pour
\(I=[a_-,a_+]\),
\[
 F_0(a_-)<0<F_0(a_+),
 \qquad
 m_I:=\inf_{a\in I}F_0'(a)>0.
\]
Si \(R\in C^1(I)\) satisfait
\begin{equation}
 \|R\|_{C^0(I)}
 <\delta_I
 :=\min\{-F_0(a_-),F_0(a_+)\},
 \qquad
 \|R'\|_{C^0(I)}<m_I,
 \label{eq:rebote-c1-bounds}
\end{equation}
alors il existe un unique \(a_R\in(a_-,a_+)\) avec \(F_R(a_R)=0\), et
ce zéro est transversal :
\begin{equation}
 F_R'(a_R)>0,
 \qquad
 \dot H_R(a_R)=\frac{a_R}{2}F_R'(a_R)>0.
 \label{eq:rebote-perturbed-transversality}
\end{equation}
Pour une famille \(R(a,\eta)\) de classe \(C^1\), le théorème des fonctions
implicites fournit une branche locale \(a_R(\eta)\) de classe \(C^1\).
\end{theorem}

\begin{proof}
L'identité \cref{eq:rebote-positive-derivative} équivaut à
\(F_0'(a_b)>0\). Par continuité, on peut choisir \(I\) avec
\(m_I>0\) et avec le changement de signe indiqué. La première borne de
\cref{eq:rebote-c1-bounds} conserve
\(F_R(a_-)<0<F_R(a_+)\) ; le théorème des valeurs intermédiaires produit un zéro.
La deuxième borne donne
\(F_R'(a)\geq m_I-\|R'\|_{C^0(I)}>0\), de sorte que \(F_R\) est strictement
croissante et que le zéro est unique. Dans chaque branche avec \(H\neq0\), dériver
\(H^2=F_R(a)\) et utiliser \(\dot a=aH\) donne
\(\dot H=(a/2)F_R'(a)\) ; la continuité étend l'égalité à \(a_R\).
La non-dégénérescence \(F_R'(a_R)>0\) est précisément l'hypothèse du
théorème des fonctions implicites.
\end{proof}

\section{Amplitude torsionnelle à partir de l'écran de Cartan}

L'exposant de la mémoire est fixé par
\cref{eq:rebote-memory-six}. L'amplitude peut être recherchée prospectivement dans
la réponse d'écran. Soit l'identité exacte

\begin{equation}
 (\beta_m^{\square})^*\Lambda_m^{\square}\beta_m^{\square}
 =\frac{56}{81}I_{Q_{\square}}.
 \label{eq:rebote-screen-energy}
\end{equation}

Une réalisation gravitationnelle complète doit construire un courant de spin

\begin{equation}
 \mathcal J_m
 =\mathcal J(\beta_m^{\square},U_m,c_s,c_g)
 \label{eq:rebote-spin-current}
\end{equation}

et une fonctionnelle positive

\begin{equation}
 s_{0,m}=\mathfrak S(\mathcal J_m,\Lambda_m^{\square})
 \label{eq:rebote-amplitude-functional}
\end{equation}

compatible avec le raffinement. La normalisation nonadique pertinente est
l'espérance traciale, car elle conserve la densité de
\cref{eq:rebote-screen-energy} ; la somme brute compte neuf copies et appartient
à une grandeur extensive distincte.

\begin{proposition}[Obligations du sélecteur d'amplitude]
Un sélecteur admissible pour \(s_0\) doit satisfaire :
\begin{enumerate}[label=\roman*)]
 \item positivité et dimensions de densité de spin au carré ;
 \item invariance de jauge sous la réduction cubique transportée ;
 \item compatibilité avec l'inversion d'arête et la retenue ;
 \item naturalité sous le raffinement \(9\)-aire ;
 \item indépendance de tout rayon de rebond utilisé comme valeur cible.
\end{enumerate}
\end{proposition}

Ces conditions transforment l'amplitude en une obligation algébrique
localisée. L'égalité énergétique d'écran fournit la forme
quadratique ; le réalisateur de Cartan doit déterminer comment cette forme se
contracte avec le courant physique.

\section{Torsion, orientation et mémoire de transport}

La métrique publie distances et cônes causaux. La connexion de Cartan
publie en outre l'orientation du transport. En HMT, chemin, retenue et
état profond constituent déjà des données de la généalogie. La correspondance

\begin{equation}
 \text{mémoire stable de voie}
 \longrightarrow
 \text{courant de spin}
 \longrightarrow
 \text{contorsion}
 \longrightarrow
 \rho_s\propto a^{-6}
 \label{eq:rebote-genealogy}
\end{equation}

offre une lecture physique de cette information orientée. La courbure
décrit la réponse de la géométrie à l'énergie ; la torsion conserve
l'histoire d'orientation qui accompagne le transport. La loi de sixième
ordre est le point où la mémoire généalogique et la densité quadratique de
spin possèdent exactement la même covariance d'échelle.

\section{Principe de Mach et auto-inertie cosmographique}

L'état local s'évalue dans un chemin compatible avec la frontière
globale. La réponse inertielle appartient donc au stabilisateur
holonomique de la totalité. Dans le secteur homogène, l'échelle \(a_n\), la
mémoire \(M_n\) et le courant torsionnel sont trois lectures du même état
généalogique enrichi :

\begin{equation}
 X_n\longmapsto
 \bigl(a_n,M_n,\mathcal J_n\bigr).
 \label{eq:rebote-mach-triple}
\end{equation}

L'interprétation machienne acquiert ainsi un contenu opératoire : la
réponse locale dépend de la classe globale de clôture et l'observateur fait
partie de la section prolongée. L'auto-inertie est la conservation de
l'état relatif sous l'holonomie qui transporte conjointement géométrie,
observateur et mémoire.

\section{Identités exactes de la dynamique torsionnelle et conditions de l'interface Einstein–Cartan}

\begin{proofstatusbox}
La loi \(M_n=(a_n/a_0)^{-6}\) est un théorème interne exact. Le théorème de
rebond \cref{eq:rebote-radius,eq:rebote-positive-derivative} est une
conséquence exacte des équations homogènes d'Einstein--Cartan
\cite{Hehl1976} avec
\(\rho_0,s_0>0\) et \(w<1\). Le transducteur
\cref{eq:rebote-transductor} conserve exactement la loi d'échelle. La
détermination de \(s_0\) à partir du courant d'écran, l'extension aux
anisotropies et la dérivation complète de l'action Einstein--Cartan sont
des interfaces physiques délimitées.
\end{proofstatusbox}

\begin{falsifierbox}
Le bloc interne est réfuté si le mode de Perron ne satisfait pas
\cref{eq:rebote-memory-six}. La réalisation d'Einstein--Cartan est réfutée
si le courant induit produit une autre puissance, une amplitude non positive ou
un signe qui élimine \cref{eq:rebote-positive-derivative}. Le rebond
homogène échoue pour les données déclarées si le rayon
\cref{eq:rebote-radius} n'existe pas, si \(H(a_b)\ne0\) ou si \(\dot H_b\le0\). La
promotion complète est rejetée si la fonctionnelle d'amplitude consulte
rétrospectivement le rayon souhaité ou perd jauge, retenue ou raffinement.
\end{falsifierbox}
